\section{Properties of the line cubes and overview of spectral setups}
\label{appendix:spectral_setups}
In \cref{tab:spwsbeams_contprop} the geometric average of the expected continuum beam from the proposal ($\theta_{prop}$) and the obtained beams ($\theta_{cont}$) are given. In \cref{tab:spwsbeams}, we provide the beam major and minor axes, along with the position angle for all observed protoclusters and spectral windows. In \cref{tab:spwsfreq}, we provide the exact range of the frequency setups for each spectral window towards all 15 protoclusters. In \cref{tab:spwsrms}, we provide the rms noise extracted for each spectral window (in units of mJy per dirty beam) as well as the noise expected from the proposed observations, scaled to the same spectral resolution and central frequency. In \cref{appendix:tab:epsilon}, we provide the $\epsilon$ factors used in the JvM correction for each spectral window.
\begin{table}[htbp!]
\centering
\caption{Proposed and actual continuum beams for all spectral windows. \label{tab:spwsbeams_contprop}}
%(using global measurements)}
\begin{tabular}{@{\extracolsep{3pt}}lccc@{}}

\hline \noalign {\smallskip}

Protocluster & & B3  & B6   \\ %& \lbol & \tdust 
%\cline{2-3}\cline{4-5}
cloud name &$\theta_{prop}$ &$\theta_{cont}$  &$\theta_{cont}$  \\
& [\arcsec]&[\arcsec]& [\arcsec] \\
\hline \noalign {\smallskip}

G008.67  & 0.67  & 0.52  & 0.66\\
G010.62  & 0.37   &0.36  & 0.47\\
G012.80  & 0.95  &1.3  & 0.88\\
G327.29  & 0.67 & 0.4   &  0.66\\
G328.25  & 0.67 & 0.51  & 0.54 \\
G337.92  & 0.51 & 0.37  & 0.54\\
G338.93  & 0.51  & 0.42   &0.53 \\
G333.60  & 0.51 & 0.46   &0.55 \\
G351.77  & 0.95 & 1.4   & 0.77\\
G353.41  & 0.95 & 1.2   & 0.79\\
W43-MM1  &0.37  & 0.41  & 0.43\\
W43-MM2  &0.37  & 0.27    &0.46 \\
W43-MM3  &0.37  & 0.35   &0.49\\
W51-E    &0.37  & 0.27   &0.3 \\
W51-IRS2 &0.37  & 0.28  & 0.47\\

\hline \noalign {\smallskip}
\end{tabular}

\end{table}



\begin{table*}[htbp!]
\centering
%\small
%    \setlength{\tabcolsep}{2pt}
%\begin{threeparttable}[c]
\caption{Obtained beams for B3 spectral windows (0-3).}
%(using global measurements)}
\label{tab:spwsbeams}
 %\begin{tabular}{lllllllllllllllllllll}
 \begin{tabular}{@{\extracolsep{3pt}}llllllllllllll@{}}

\hline \noalign {\smallskip}

Protocluster &  \multicolumn{3}{c}{spw0 }  & \multicolumn{3}{c}{spw1}  & \multicolumn{3}{c}{spw2 } & \multicolumn{3}{c}{spw3}  \\ %& \lbol & \tdust 
\cline{2-4}\cline{5-7}\cline{8-10}\cline{11-13}
cloud name&  $\theta_{maj}$ & $\theta_{min}$ &PA&  $\theta_{maj}$ & $\theta_{min}$ &PA& $\theta_{maj}$ & $\theta_{min}$ & PA& $\theta_{maj}$ & $\theta_{min}$ &PA   \\
&(\arcsec) &(\arcsec)& &(\arcsec) &(\arcsec)& & (\arcsec) &(\arcsec)& & (\arcsec) &(\arcsec)&   \\
\hline \noalign {\smallskip}

G008.67  &    0.95&0.72&72&0.55&0.43&69&0.50&0.40&65&0.49&0.38&61 \\
G010.62  &    0.53&0.39&102&0.46&0.38&107&0.37&0.28&94&0.37&0.29&107 \\
G012.80  &   2.23&1.83&88&1.85&1.65&80&1.73&1.52&84&1.70&1.43&87 \\
G327.29  &   0.66&0.57&53&0.45&0.39&69&0.41&0.35&69&0.41&0.35&68 \\
G328.25  &   0.97&0.90&71&0.69&0.51&106&0.63&0.48&113&0.61&0.45&112 \\
G337.92  &   0.67&0.60&-72&0.49&0.42&73&0.44&0.39&73&0.43&0.36&75 \\
G338.93  &   0.59&0.54&-62&0.45&0.42&40&0.39&0.35&88&0.38&0.35&84 \\
G333.60  &   0.73&0.70&13&0.50&0.47&45&0.45&0.44&25&0.45&0.42&37 \\
G351.77  &   2.20&1.88&91&1.85&1.65&88&1.69&1.51&89&1.72&1.49&100 \\
G353.41  &   2.12&1.76&77&1.83&1.54&69&1.67&1.50&76&1.62&1.47&78 \\
W43-MM1  &    0.88&0.48&99&0.59&0.37&107&0.54&0.33&107&0.53&0.33&108 \\
W43-MM2  &   0.41&0.32&106&0.34&0.26&106&0.30&0.24&110&0.32&0.24&124 \\
W43-MM3  &  0.58&0.42&97&0.47&0.33&99&0.42&0.29&98&0.41&0.29&99 \\
W51-E  &   0.38&0.35&106&0.36&0.33&108&0.29&0.27&83&0.28&0.26&85 \\
W51-IRS2  &   0.39&0.36&111&0.33&0.30&131&0.29&0.27&130&0.28&0.27&128 \\
%Average  &  & & & & 1.52  &   &  & 1.15  &  &  &  1.02  &  &  &  1.01 \\

\hline \noalign {\smallskip}
\end{tabular}
%\begin{tablenotes}

%
%\end{tablenotes}
%\end{threeparttable}
\end{table*}
%\pagestyle{empty}
%\begin{landscape}
\begin{table*}[htbp!]
\ContinuedFloat
\centering
%\small
%    \setlength{\tabcolsep}{2pt}
%\begin{threeparttable}[c]
\caption{Continued: Obtained beams for B6 spectral windows (0-3).}
%\label{tab:spwsbeams_2}
%(using global measurements)}
 \begin{tabular}{@{\extracolsep{3pt}}llllllllllllll@{}}

\hline \noalign {\smallskip}

Protocluster &  \multicolumn{3}{c}{spw0 }  & \multicolumn{3}{c}{spw1}  & \multicolumn{3}{c}{spw2 } & \multicolumn{3}{c}{spw3}  \\ %& \lbol & \tdust 
\cline{2-4}\cline{5-7}\cline{8-10}\cline{11-13}
cloud name&  $\theta_{maj}$ & $\theta_{min}$ &PA&  $\theta_{maj}$ & $\theta_{min}$ &PA& $\theta_{maj}$ & $\theta_{min}$ & PA& $\theta_{maj}$ & $\theta_{min}$ &PA   \\
&(\arcsec) &(\arcsec)& &(\arcsec) &(\arcsec)& & (\arcsec) &(\arcsec)& & (\arcsec) &(\arcsec)&   \\
\hline \noalign {\smallskip}
G008.67  &  0.88&0.72&98&0.88&0.72&98&0.88&0.72&96&0.88&0.71&99 \\
G010.62  &  0.65&0.51&106&0.68&0.53&-72&0.64&0.50&106&0.66&0.52&-76 \\
G012.80  &  1.30&0.89&76&1.29&0.88&76&1.29&0.86&76&1.29&0.87&77 \\
G327.29  &  0.82&0.76&121&0.82&0.75&124&0.82&0.74&123&0.82&0.74&125 \\
G328.25  &  0.74&0.59&-13&0.74&0.58&-14&0.75&0.59&-14&0.74&0.58&-14 \\
G337.92  &  0.81&0.66&129&0.80&0.66&129&0.79&0.66&128&0.80&0.66&129 \\
G338.93  &  0.77&0.69&80&0.77&0.68&81&0.74&0.65&81&0.75&0.65&83 \\
G333.60  &  0.75&0.69&-35&0.75&0.68&-36&0.74&0.68&-37&0.74&0.68&-38 \\
G351.77  &  1.08&0.84&88&1.08&0.83&87&1.07&0.82&88&1.08&0.82&88 \\
G353.41  &  1.13&0.83&86&1.13&0.83&85&1.11&0.83&86&1.12&0.83&86 \\
W43-MM1  &  0.66&0.48&-81&0.65&0.47&100&0.67&0.49&103&0.65&0.49&105 \\
W43-MM2  &  0.63&0.52&100&0.62&0.51&-84&0.68&0.54&100&0.64&0.50&96 \\
W43-MM3  &  0.66&0.57&86&0.66&0.57&86&0.66&0.57&86&0.66&0.57&86 \\
W51-E  & 0.46&0.35&29&0.46&0.35&30&0.57&0.49&38&0.46&0.35&30 \\
W51-IRS2  &  0.64&0.57&-18&0.65&0.59&-23&0.65&0.58&-21&0.63&0.56&-24 \\
\hline \noalign {\smallskip}
\end{tabular}

%\end{threeparttable}
\end{table*}

\begin{table*}
\ContinuedFloat
\centering
%\small
%    \setlength{\tabcolsep}{2pt}
%\begin{threeparttable}[c]
\caption{Continued: Obtained beams for B6 spectral windows (4-7).}
%\label{tab:spwsbeams_3}
%(using global measurements)}
 \begin{tabular}{@{\extracolsep{3pt}}llllllllllllll@{}}

\hline \noalign {\smallskip}

Protocluster &  \multicolumn{3}{c}{spw4 }  & \multicolumn{3}{c}{spw5}  & \multicolumn{3}{c}{spw6 } & \multicolumn{3}{c}{spw7}  \\ %& \lbol & \tdust 
\cline{2-4}\cline{5-7}\cline{8-10}\cline{11-13}
cloud name&  $\theta_{maj}$ & $\theta_{min}$ &PA&  $\theta_{maj}$ & $\theta_{min}$ &PA& $\theta_{maj}$ & $\theta_{min}$ & PA& $\theta_{maj}$ & $\theta_{min}$ &PA   \\
&(\arcsec) &(\arcsec)& &(\arcsec) &(\arcsec)& & (\arcsec) &(\arcsec)& & (\arcsec) &(\arcsec)&   \\
\hline \noalign {\smallskip}
G008.67  &   0.88&0.71&99&0.84&0.68&96&0.83&0.67&96&0.73&0.61&97 \\
G010.62  &   0.67&0.52&112&0.61&0.48&-73&0.61&0.47&106&0.58&0.43&110 \\
G012.80  &   1.29&0.87&77&1.21&0.83&76&1.21&0.82&76&1.13&0.72&76 \\
G327.29  &   0.82&0.74&125&0.77&0.70&128&0.77&0.69&127&0.69&0.62&-35 \\
G328.25  &   0.74&0.58&-14&0.69&0.54&-13&0.69&0.54&-13&0.61&0.45&-13 \\
G337.92  &   0.80&0.66&129&0.75&0.61&127&0.74&0.61&127&0.65&0.51&127 \\
G338.93  &   0.75&0.66&82&0.71&0.63&84&0.71&0.63&83&0.60&0.55&91 \\
G333.60  &   0.74&0.68&-36&0.69&0.62&-34&0.68&0.62&-34&0.57&0.52&-23 \\
G351.77  &   1.07&0.82&88&1.01&0.77&88&1.00&0.77&87&0.91&0.68&87 \\
G353.41  &   1.12&0.82&84&1.06&0.77&84&1.05&0.76&86&0.96&0.67&85 \\
W43-MM1  &   0.69&0.50&99&0.59&0.43&100&0.59&0.43&100&0.50&0.34&103 \\
W43-MM2  &   0.68&0.53&96&0.61&0.51&110&0.61&0.51&110&0.56&0.45&112 \\
W43-MM3  &   0.66&0.57&85&0.62&0.54&86&0.62&0.53&85&0.55&0.46&88 \\
W51-E  &   0.57&0.47&39&0.52&0.41&34&0.42&0.32&30&0.34&0.27&25 \\
W51-IRS2   &  0.63&0.56&-23&0.59&0.52&-23&0.59&0.52&-25&0.54&0.47&-26 \\
%Average  &  & & & & 1.27  &   &  & 1.19  &  &  &  1.16  &  &  &  1.02 \\
\hline \noalign {\smallskip}
\end{tabular}
\end{table*}
\begin{table*}[htbp!]
\centering
%\small
%    \setlength{\tabcolsep}{2pt}
%\begin{threeparttable}[c]
\caption{Minimum and maximum frequency ranges [GHz] for the Spectral window setups for B3 spws (0-3).}
%(using global measurements)}
%\label{tab:spwsfreq_1}
\label{tab:spwsfreq}

\begin{tabular}{@{\extracolsep{6pt}}lcccccccc@{}}
 %\begin{tabular}{lrlcrlrlrlrlr}
\hline \noalign {\smallskip}
Protocluster  & \multicolumn{2}{c}{spw0} & \multicolumn{2}{c}{spw1} & \multicolumn{2}{c}{spw2}& \multicolumn{2}{c}{spw3} \\ 
\cline{2-3}\cline{4-5}\cline{6-7}\cline{8-9}
 %\cpartlineleft{1,1em}
%\cline{5-6} \cline{7-8}   
cloud name\tnote{1}  & min & max  & min & max  & min & max  & min & max  \\
\hline \noalign {\smallskip}

G008.67 & 93.103&93.220&91.720&92.656&102.119&103.056&104.519&105.455 \\
G010.62 & 93.116&93.233&91.733&92.669&102.133&103.069&104.533&105.469 \\
G012.80 & 93.103&93.220&91.720&92.657&102.119&103.056&104.519&105.455 \\
G327.29 & 93.129&93.246&91.746&92.682&102.147&103.084&104.547&105.484 \\
G328.25 & 93.128&93.245&91.745&92.682&102.147&103.083&104.547&105.483 \\
G337.92 & 93.127&93.244&91.744&92.681&102.146&103.082&104.546&105.482 \\
G338.93 & 93.134&93.251&91.751&92.688&102.153&103.089&104.554&105.490 \\
G333.60 & 93.130&93.247&91.746&92.683&102.148&103.085&104.548&105.485 \\
G351.77 & 93.116&93.233&91.733&92.669&102.133&103.069&104.533&105.469 \\
G353.41 & 93.120&93.237&91.737&92.674&102.138&103.074&104.538&105.474 \\
W43-MM1 & 93.085&93.202&91.702&92.639&102.099&103.035&104.498&105.434 \\
W43-MM2 & 93.085&93.202&91.702&92.638&102.099&103.035&104.498&105.435 \\
W43-MM3 & 93.085&93.202&91.702&92.639&102.099&103.035&104.498&105.434 \\
W51-E & 93.098&93.215&91.715&92.651&102.113&103.049&104.513&105.449 \\
W51-IRS2 & 93.098&93.215&91.715&92.651&102.113&103.050&104.513&105.449 \\

\hline \noalign {\smallskip}
\end{tabular}
%\begin{tablenotes}

%
%\end{tablenotes}
%\end{threeparttable}
\end{table*}
%\pagestyle{empty}
%\begin{landscape}
\begin{table*}[htbp!]
\ContinuedFloat
\centering
%\small
%    \setlength{\tabcolsep}{2pt}
%\begin{threeparttable}[c]
\caption{Continued: minimum and maximum frequency ranges [GHz] for the Spectral window setups for for B6 spws (0-3).}
%\label{tab:spwsfreq_2}
%(using global measurements)}
\begin{tabular}{@{\extracolsep{6pt}}lcccccccc@{}}
 %\begin{tabular}{lrlcrlrlrlrlr}
\hline \noalign {\smallskip}
Protocluster  & \multicolumn{2}{c}{spw0} & \multicolumn{2}{c}{spw1} & \multicolumn{2}{c}{spw2}& \multicolumn{2}{c}{spw3} \\ 
\cline{2-3}\cline{4-5}\cline{6-7}\cline{8-9}
 %\cpartlineleft{1,1em}
%\cline{5-6} \cline{7-8}   
cloud name\tnote{1}  & min & max  & min & max  & min & max  & min & max  \\
\hline \noalign {\smallskip}
G008.67 & 216.056&216.290&217.006&217.240&219.859&219.976&218.086&218.320 \\
G010.62 & 216.084&216.318&217.034&217.268&219.888&220.005&218.114&218.348 \\
G012.80 & 216.056&216.290&217.006&217.240&219.860&219.976&218.086&218.320 \\
G327.29 & 216.115&216.350&217.066&217.299&219.920&220.036&218.146&218.380 \\
G328.25 & 216.114&216.348&217.064&217.298&219.918&220.035&218.144&218.378 \\
G337.92 & 216.112&216.346&217.062&217.296&219.916&220.033&218.142&218.376 \\
G338.93 & 216.128&216.362&217.078&217.312&219.932&220.049&218.158&218.392 \\
G333.60 & 216.117&216.351&217.067&217.301&219.921&220.038&218.147&218.381 \\
G351.77 & 216.085&216.319&217.035&217.269&219.889&220.006&218.115&218.349 \\
G353.41 & 216.095&216.329&217.045&217.279&219.899&220.016&218.125&218.360 \\
W43-MM1 & 216.014&216.249&216.964&217.198&219.817&219.934&218.044&218.278 \\
W43-MM2 & 216.013&216.247&216.963&217.197&219.816&219.932&218.042&218.276 \\
W43-MM3 & 216.013&216.247&216.963&217.197&219.816&219.932&218.042&218.277 \\
W51-E & 216.043&216.277&216.993&217.227&219.846&219.963&218.073&218.307 \\
W51-IRS2 & 216.043&216.277&216.993&217.227&219.846&219.963&218.073&218.307 \\
\hline \noalign {\smallskip}
\end{tabular}
%\begin{tablenotes}

%\end{tablenotes}
%\end{threeparttable}
\end{table*}



\begin{table*}[htbp!]
\ContinuedFloat
\centering
\caption{Continued: minimum and maximum frequency ranges [GHz] for the Spectral window setups for B6 spws (4-7).}
%\label{tab:spwsfreq_3}
\begin{tabular}{@{\extracolsep{6pt}}lcccccccc@{}}
\hline \noalign {\smallskip}
Protocluster  & \multicolumn{2}{c}{spw4} & \multicolumn{2}{c}{spw5} & \multicolumn{2}{c}{spw6}& \multicolumn{2}{c}{spw7} \\ 
\cline{2-3}\cline{4-5}\cline{6-7}\cline{8-9}
  
cloud name\tnote{1}  & min & max  & min & max  & min & max  & min & max  \\
\hline \noalign {\smallskip}
G008.67 & 219.474&219.591&230.268&230.735&231.017&231.485&231.480&233.353 \\
G010.62 & 219.503&219.620&230.298&230.765&231.047&231.516&231.508&233.381 \\
G012.80 & 219.474&219.591&230.269&230.735&231.017&231.486&231.480&233.353 \\
G327.29 & 219.534&219.651&230.331&230.798&231.081&231.549&231.541&233.414 \\
G328.25 & 219.533&219.650&230.330&230.797&231.079&231.547&231.540&233.413 \\
G337.92 & 219.531&219.648&230.328&230.794&231.077&231.545&231.538&233.361 \\
G338.93 & 219.547&219.664&230.345&230.811&231.094&231.562&231.554&233.426 \\
G333.60 & 219.536&219.653&230.333&230.800&231.082&231.550&231.543&233.415 \\
G351.77 & 219.504&219.621&230.299&230.766&231.048&231.517&231.510&233.383 \\
G353.41 & 219.514&219.631&230.310&230.777&231.059&231.527&231.521&233.393 \\
W43-MM1 & 219.432&219.549&230.224&230.690&230.973&231.441&232.490&234.362 \\
W43-MM2 & 219.431&219.548&230.222&230.689&230.971&231.439&231.436&233.308 \\
W43-MM3 & 219.431&219.547&230.222&230.688&230.971&231.439&231.435&233.308 \\
W51-E & 219.461&219.578&230.255&230.721&231.004&231.472&231.467&233.339 \\
W51-IRS2 & 219.461&219.578&230.255&230.721&231.003&231.472&231.467&233.339 \\
\hline \noalign {\smallskip}
\end{tabular}
\end{table*}
\begin{table*}[htbp!]
\centering
%\small
%    \setlength{\tabcolsep}{2pt}
%\begin{threeparttable}[c]
\caption{Spectral window noise estimates for B3 spws (0-3).}
%(using global measurements)}
\label{tab:spwsrms}
 %\begin{tabular}{lllllllllllllllllllll}
\begin{tabular}{@{\extracolsep{4pt}}lllllllllllll@{}}

\hline \noalign {\smallskip}
Protocluster  & \multicolumn{3}{c}{spw0}  & \multicolumn{3}{c}{spw1}  & \multicolumn{3}{c}{spw2}
    & \multicolumn{3}{c}{spw3}  \\ %& \lbol & \tdust 
\cline{2-4}\cline{5-7}\cline{8-10}\cline{11-13}
cloud name\tnote{1}  & $\sigma_{MAD}$ & $\sigma_{req}$ & $\frac{\sigma_{MAD}}{\sigma_{req}}$  & $\sigma_{MAD}$ & $\sigma_{req}$ & $\frac{\sigma_{MAD}}{\sigma_{req}}$  & $\sigma_{MAD}$ & $\sigma_{req}$ & $\frac{\sigma_{MAD}}{\sigma_{req}}$  & $\sigma_{MAD}$ & $\sigma_{req}$ & $\frac{\sigma_{MAD}}{\sigma_{req}}$\\
& \multicolumn{2}{c}{[mJybeam$^{-1}$]}& & \multicolumn{2}{c}{[mJybeam$^{-1}$]}& &  \multicolumn{2}{c}{[mJybeam$^{-1}$]}& & \multicolumn{2}{c}{[mJybeam$^{-1}$]}\\
\hline \noalign {\smallskip}
G008.67  &  7.39 & 14.84 & 0.50 & 2.78 & 5.22 & 0.53 & 3.28 & 5.51 & 0.60 & 3.59 & 5.57 & 0.64 \\
G010.62  &  3.52 & 5.04 & 0.70 & 1.51 & 1.77 & 0.85 & 1.53 & 1.87 & 0.82 & 1.32 & 1.89 & 0.70 \\
G012.80  &  11.18 & 28.36 & 0.39 & 3.95 & 9.97 & 0.40 & 4.82 & 10.52 & 0.46 & 5.10 & 10.64 & 0.48 \\
G327.29  &  6.23 & 14.92 & 0.42 & 1.26 & 5.25 & 0.24 & 2.86 & 5.54 & 0.52 & 2.88 & 5.60 & 0.52 \\
G328.25  &  8.16 & 14.83 & 0.55 & 5.38 & 5.22 & 1.03 & 3.52 & 5.50 & 0.64 & 3.76 & 5.57 & 0.67 \\
G337.92  &  3.63 & 9.93 & 0.37 & 2.25 & 3.49 & 0.64 & 2.52 & 3.68 & 0.68 & 1.60 & 3.73 & 0.43 \\
G338.93  &  4.67 & 9.89 & 0.47 & 2.15 & 3.48 & 0.62 & 2.47 & 3.67 & 0.67 & 2.66 & 3.71 & 0.72 \\
G333.60  &  4.97 & 9.84 & 0.50 & 1.84 & 3.46 & 0.53 & 2.02 & 3.65 & 0.55 & 2.25 & 3.69 & 0.61 \\
G351.77  &  10.84 & 28.50 & 0.38 & 3.78 & 10.02 & 0.38 & 4.39 & 10.58 & 0.41 & 5.69 & 10.70 & 0.53 \\
G353.41  &  10.75 & 28.46 & 0.38 & 3.91 & 10.01 & 0.39 & 4.45 & 10.56 & 0.42 & 4.95 & 10.68 & 0.46 \\
W43-MM1  &  2.60 & 4.99 & 0.52 & 0.88 & 1.75 & 0.50 & 1.07 & 1.85 & 0.58 & 0.96 & 1.87 & 0.51 \\
W43-MM2  &  2.68 & 4.99 & 0.54 & 1.09 & 1.75 & 0.62 & 1.26 & 1.85 & 0.68 & 1.38 & 1.87 & 0.74 \\
W43-MM3  &  3.14 & 4.99 & 0.63 & 1.15 & 1.75 & 0.66 & 1.29 & 1.85 & 0.70 & 1.42 & 1.87 & 0.76 \\
W51-E  &  2.23 & 4.95 & 0.45 & 0.92 & 1.74 & 0.53 & 1.06 & 1.84 & 0.58 & 1.16 & 1.86 & 0.62 \\
W51-IRS2  &  2.58 & 4.95 & 0.52 & 0.93 & 1.74 & 0.54 & 1.24 & 1.84 & 0.67 & 1.29 & 1.86 & 0.69 \\
\hline \noalign {\smallskip}
\end{tabular}
\begin{tablenotes}
\item $\sigma_{MAD}$ is given in units of mJy per dirty beam as discussed in \cref{sec:jvm}, while $\sigma_{req}$ is the noise expected from the proposed ALMA observations.
\end{tablenotes}
%\end{threeparttable}
\end{table*}
%\pagestyle{empty}
%\begin{landscape}
\begin{table*}[htbp!]
\ContinuedFloat
\centering
%\small
%    \setlength{\tabcolsep}{2pt}
%\begin{threeparttable}[c]
\caption{Spectral window noise estimates for B6 spws (0-3).}
%\label{tab:spwsrms_2}
%(using global measurements)}
\begin{tabular}{@{\extracolsep{4pt}}lllllllllllll@{}}

\hline \noalign {\smallskip}
Protocluster  & \multicolumn{3}{c}{spw0}  & \multicolumn{3}{c}{spw1}  & \multicolumn{3}{c}{spw2}
    & \multicolumn{3}{c}{spw3} \\ %& \lbol & \tdust 
\cline{2-4}\cline{5-7}\cline{8-10}\cline{11-13}
cloud name\tnote{1}  & $\sigma_{MAD}$ & $\sigma_{req}$ & $\frac{\sigma_{MAD}}{\sigma_{req}}$  & $\sigma_{MAD}$ & $\sigma_{req}$ & $\frac{\sigma_{MAD}}{\sigma_{req}}$  & $\sigma_{MAD}$ & $\sigma_{req}$ & $\frac{\sigma_{MAD}}{\sigma_{req}}$  & $\sigma_{MAD}$ & $\sigma_{req}$ & $\frac{\sigma_{MAD}}{\sigma_{req}}$\\
& \multicolumn{2}{c}{[mJybeam$^{-1}$]}& & \multicolumn{2}{c}{[mJybeam$^{-1}$]}& &  \multicolumn{2}{c}{[mJybeam$^{-1}$]}& & \multicolumn{2}{c}{[mJybeam$^{-1}$]}\\
\hline \noalign {\smallskip}
G008.67  &  8.15 & 33.93 & 0.24 & 7.20 & 22.40 & 0.32 & 10.01 & 22.55 & 0.44 & 7.54 & 33.93 & 0.22 \\
G010.62  &  2.38 & 11.69 & 0.20 & 1.98 & 7.72 & 0.26 & 2.65 & 7.77 & 0.34 & 2.21 & 11.69 & 0.19 \\
G012.80  &  12.40 & 53.77 & 0.23 & 11.27 & 35.50 & 0.32 & 16.30 & 35.73 & 0.46 & 11.52 & 53.77 & 0.21 \\
G327.29  &  9.78 & 37.01 & 0.26 & 8.49 & 24.44 & 0.35 & 10.15 & 24.59 & 0.41 & 8.43 & 37.01 & 0.23 \\
G328.25  &  15.10 & 36.99 & 0.41 & 13.58 & 24.43 & 0.56 & 17.84 & 24.58 & 0.73 & 14.29 & 36.99 & 0.39 \\
G337.92  &  3.84 & 22.77 & 0.17 & 2.98 & 15.04 & 0.20 & 4.36 & 15.13 & 0.29 & 3.39 & 22.77 & 0.15 \\
G338.93  &  3.67 & 22.09 & 0.17 & 3.19 & 14.59 & 0.22 & 4.45 & 14.68 & 0.30 & 3.48 & 22.09 & 0.16 \\
G333.60  &  3.18 & 22.84 & 0.14 & 2.59 & 15.08 & 0.17 & 3.74 & 15.18 & 0.25 & 2.92 & 22.84 & 0.13 \\
G351.77  &  13.22 & 56.26 & 0.23 & 11.60 & 37.15 & 0.31 & 15.79 & 37.39 & 0.42 & 11.77 & 56.20 & 0.21 \\
G353.41  &  13.94 & 56.75 & 0.25 & 12.59 & 37.47 & 0.34 & 16.44 & 37.71 & 0.44 & 12.60 & 56.75 & 0.22 \\
W43-MM1  &  2.18 & 11.77 & 0.19 & 1.94 & 7.77 & 0.25 & 2.99 & 7.82 & 0.38 & 2.04 & 11.77 & 0.17 \\
W43-MM2  &  1.95 & 11.77 & 0.17 & 1.63 & 7.77 & 0.21 & 2.52 & 7.82 & 0.32 & 1.87 & 11.77 & 0.16 \\
W43-MM3  &  2.08 & 11.77 & 0.18 & 1.86 & 7.77 & 0.24 & 2.46 & 7.82 & 0.31 & 1.86 & 11.77 & 0.16 \\
W51-E  &  2.10 & 12.00 & 0.17 & 1.85 & 7.96 & 0.23 & 2.11 & 8.01 & 0.26 & 2.00 & 12.05 & 0.17 \\
W51-IRS2  &  2.30 & 12.05 & 0.19 & 2.07 & 7.96 & 0.26 & 2.82 & 8.01 & 0.35 & 2.11 & 12.05 & 0.17 \\


\hline \noalign {\smallskip}
\end{tabular}
\begin{tablenotes}
\item $\sigma_{MAD}$ is given in units of mJy per dirty beam as discussed in \cref{sec:jvm}, while $\sigma_{req}$ is the noise expected from the proposed ALMA observations.
\end{tablenotes}
%\end{threeparttable}
\end{table*}


\begin{table*}[htbp!]
\ContinuedFloat
\centering
%\small
%    \setlength{\tabcolsep}{2pt}
%\begin{threeparttable}[c]
\caption{Spectral window noise estimates for B6 spws (4-7).}
%\label{tab:spwsrms_3}
%(using global measurements)}
\begin{tabular}{@{\extracolsep{4pt}}lllllllllllll@{}}

\hline \noalign {\smallskip}
Protocluster  & \multicolumn{3}{c}{spw4}  & \multicolumn{3}{c}{spw5}  & \multicolumn{3}{c}{spw6}
    & \multicolumn{3}{c}{spw7}  \\ %& \lbol & \tdust 
\cline{2-4}\cline{5-7}\cline{8-10}\cline{11-13}
cloud name\tnote{1}  & $\sigma_{MAD}$ & $\sigma_{req}$ & $\frac{\sigma_{MAD}}{\sigma_{req}}$  & $\sigma_{MAD}$ & $\sigma_{req}$ & $\frac{\sigma_{MAD}}{\sigma_{req}}$  & $\sigma_{MAD}$ & $\sigma_{req}$ & $\frac{\sigma_{MAD}}{\sigma_{req}}$  & $\sigma_{MAD}$ & $\sigma_{req}$ & $\frac{\sigma_{MAD}}{\sigma_{req}}$\\
& \multicolumn{2}{c}{[mJybeam$^{-1}$]}& & \multicolumn{2}{c}{[mJybeam$^{-1}$]}& &  \multicolumn{2}{c}{[mJybeam$^{-1}$]}& & \multicolumn{2}{c}{[mJybeam$^{-1}$]}\\
\hline \noalign {\smallskip}
G008.67  &  15.89 & 33.93 & 0.47 & 7.67 & 12.46 & 0.62 & 10.09 & 17.58 & 0.57 & 6.43 & 11.59 & 0.56 \\
G010.62  &  5.15 & 11.69 & 0.44 & 2.61 & 4.29 & 0.61 & 3.21 & 6.06 & 0.53 & 1.75 & 3.99 & 0.44 \\
G012.80  &  21.63 & 53.77 & 0.40 & 10.76 & 19.74 & 0.55 & 11.90 & 27.87 & 0.43 & 7.01 & 18.37 & 0.38 \\
G327.29  &  16.42 & 37.01 & 0.44 & 8.14 & 13.59 & 0.60 & 10.45 & 19.18 & 0.54 & 6.41 & 12.65 & 0.51 \\
G328.25  &  32.64 & 36.99 & 0.88 & 16.99 & 13.58 & 1.25 & 22.07 & 19.17 & 1.15 & 22.71 & 12.64 & 1.80 \\
G337.92  &  8.48 & 22.77 & 0.37 & 3.67 & 8.36 & 0.44 & 5.05 & 11.80 & 0.43 & 2.80 & 7.78 & 0.36 \\
G338.93  &  8.45 & 22.09 & 0.38 & 3.38 & 8.11 & 0.42 & 5.11 & 11.45 & 0.45 & 2.74 & 7.55 & 0.36 \\
G333.60  &  7.50 & 22.84 & 0.33 & 3.04 & 8.39 & 0.36 & 4.57 & 11.84 & 0.39 & 2.45 & 7.81 & 0.31 \\
G351.77  &  20.76 & 56.26 & 0.37 & 12.15 & 20.66 & 0.59 & 13.36 & 29.16 & 0.46 & 8.49 & 19.23 & 0.44 \\
G353.41  &  22.05 & 56.70 & 0.39 & 12.74 & 20.84 & 0.61 & 13.58 & 29.41 & 0.46 & 8.77 & 19.39 & 0.45 \\
W43-MM1  &  4.69 & 11.77 & 0.40 & 2.43 & 4.32 & 0.56 & 3.05 & 6.10 & 0.50 & 2.09 & 4.02 & 0.52 \\
W43-MM2  &  4.25 & 11.77 & 0.36 & 1.58 & 4.32 & 0.37 & 2.56 & 6.10 & 0.42 & 1.63 & 4.02 & 0.40 \\
W43-MM3  &  4.64 & 11.77 & 0.39 & 2.27 & 4.32 & 0.53 & 2.77 & 6.10 & 0.45 & 1.62 & 4.02 & 0.40 \\
W51-E  &  4.71 & 12.05 & 0.39 & 2.78 & 4.43 & 0.63 & 3.57 & 6.25 & 0.57 & 2.45 & 4.12 & 0.59 \\
W51-IRS2  &  4.41 & 12.05 & 0.37 & 2.05 & 4.43 & 0.46 & 2.90 & 6.25 & 0.46 & 1.55 & 4.12 & 0.38 \\

\hline \noalign {\smallskip}
\end{tabular}
\begin{tablenotes}
\item $\sigma_{MAD}$ is given in units of mJy per dirty beam as discussed in \cref{sec:jvm}, while $\sigma_{req}$ is the noise expected from the proposed ALMA observations.

\end{tablenotes}
%\end{threeparttable}
\end{table*}
\begin{table*}[htbp!]
\centering
\caption{\label{appendix:tab:epsilon}The $\epsilon$ values used for the JvM correction for B3 spws (0-3) and B6 spws (0-7). }
\begin{tabular}{@{\extracolsep{6pt}}lllllllllllll@{}}
 \hline \noalign {\smallskip}
 Protocluster  & \multicolumn{4}{c}{B3 }& \multicolumn{8}{c}{B6} \\ %& \lbol & \tdust 
 \cline{2-5} \cline{6-13}
cloud name & spw0& spw1  & spw2 & spw3 & spw0  & spw1  & spw2 & spw3 & spw4 & spw5 & spw6 & spw7 \\
\hline \noalign {\smallskip}
G008.67 & 0.37&0.49&0.47&0.47 & 0.67&0.67&0.67&0.67&0.66&0.67&0.67&0.64 \\
G010.62 & 0.32&0.47&0.94&0.46 & 0.61&0.61&0.61&0.61&0.61&0.6&0.6&0.87 \\
G012.80 & 0.66&0.56&0.58&0.59 & 0.84&0.84&0.85&0.84&0.84&0.84&0.84&0.81 \\
G327.29 & 0.23&0.88&0.38&0.93 & 0.7&0.7&0.69&0.7&0.7&0.68&0.69&0.67 \\
G328.25 & 0.32&0.33&0.52&0.52 & 0.6&0.6&0.59&0.6&0.6&0.6&0.6&0.61 \\
G337.92 & 0.4&0.31&0.32&0.48 & 0.59&0.59&0.58&0.59&0.59&0.58&0.58&0.52 \\
G338.93 & 0.31&0.38&0.25&0.25 & 0.58&0.58&0.58&0.57&0.58&0.58&0.57&0.5 \\
G333.60 & 0.31&0.48&0.5&0.46 & 0.54&0.54&0.54&0.54&0.54&0.52&0.52&0.49 \\
G351.77 & 0.67&0.58&0.58&0.6 & 0.79&0.81&0.79&0.74&0.79&0.74&0.76&0.7 \\
G353.41 & 0.67&0.56&0.6&0.58 & 0.78&0.79&0.74&0.79&0.75&0.76&0.79&0.71 \\
W43-MM1 & 0.47&0.95&0.37&0.4 & 0.6&0.6&0.62&0.62&0.62&0.61&0.61&0.61 \\
W43-MM2 & 0.26&0.94&0.97&0.95 & 0.65&0.62&0.66&0.65&0.65&0.64&0.64&0.65 \\
W43-MM3 & 0.32&0.85&1.01&1.0 & 0.59&0.59&0.59&0.58&0.58&0.58&0.57&0.58 \\
W51-E & 0.42&0.26&0.99&0.98 & 0.48&0.48&0.43&0.48&0.47&0.4&0.47&0.48 \\
W51-IRS2 & 0.41&0.96&0.99&0.98 & 0.65&0.65&0.65&0.65&0.65&0.65&0.65&0.98 \\
\hline \noalign {\smallskip}
\end{tabular}
\end{table*}
\section{Known data issues}

Towards several regions that contain bright hot cores, there can be a slight overestimation of the continuum startmodel at the positions of the hot cores, resulting in localized negative features in some of the residuals. This overestimate is likely due to rich line emission (line forests) across the spectrum at those positions. Often the negative residuals are < 10\% of the flux in the restored image, but in some cases with faint (real) continuum and bright line emission, they are larger. We also note that in some instances the sidelobes can remain particularly bright in the final imaged cubes (e.g., G327 in B6). 